\section{Teorema Rayleigh-Ritz: Fondasi Matematis Prinsip Variasi}\label{appendix-d}

\hypertarget{filosofi-kecenderungan-alam-meminimalkan-energi}{%
\subsection{Filosofi: Kecenderungan Alam Meminimalkan Energi}\label{filosofi-kecenderungan-alam-meminimalkan-energi}}

Dalam fisika, sistem selalu cenderung menuju kondisi dengan energi terendah (\emph{Ground State}). \textbf{Prinsip Variasi} adalah alat matematika yang memungkinkan kita ``menebak'' kondisi tersebut dengan jaminan bahwa tebakan kita tidak akan pernah lebih rendah dari kenyataan alam semesta.

\begin{center}\rule{0.5\linewidth}{0.5pt}\end{center}

\hypertarget{kalkulus-variasi-mencari-titik-stasioner-energi}{%
\subsection{Kalkulus Variasi: Mencari Titik Stasioner Energi}\label{kalkulus-variasi-mencari-titik-stasioner-energi}}

Bayangkan energi sistem bukan sebagai angka, melainkan sebagai sebuah \textbf{fungsional} \(E[\psi]\) yang bergantung pada bentuk fungsi gelombang (atau state) \(|\psi\rangle\).

\hypertarget{a.-definisi-fungsional-energi}{%
\subsubsection{Definisi Fungsional Energi}\label{a.-definisi-fungsional-energi}}

Energi rata-rata dari sebuah sistem didefinisikan sebagai hasil bagi Rayleigh: \[ E(\psi) = \frac{\langle \psi | \hat{H} | \psi \rangle}{\langle \psi | \psi \rangle} \qquad (1) \]

\begin{quote}
\textbf{Visualisasi (1): Jembatan Logika} Mengapa kita membagi dengan \(\langle \psi | \psi \rangle\)? Karena kita ingin memastikan nilai energi tidak bergantung pada ``panjang'' vektor \(|\psi\rangle\) (normalisasi). Jika \(|\psi\rangle\) sudah ternormalisasi (\(\langle \psi | \psi \rangle = 1\)), maka \(E(\psi) = \langle \psi | \hat{H} | \psi \rangle\).
\end{quote}

\hypertarget{b.-variasi-fungsional-delta-e-0}{%
\subsubsection{\texorpdfstring{Variasi Fungsional (\(\delta E = 0\))}{Variasi Fungsional (\textbackslash delta E = 0)}}\label{b.-variasi-fungsional-delta-e-0}}

Untuk mencari energi minimum, kita lakukan variasi kecil terhadap \(|\psi\rangle\) dan menetapkan perubahannya menjadi nol. \[ \delta E = \delta \left[ \frac{\langle \psi | \hat{H} | \psi \rangle}{\langle \psi | \psi \rangle} \right] = 0 \qquad (2) \]

\begin{quote}
\textbf{Visualisasi (2): Perhitungan Linear (Uraian Variasi)} Menggunakan aturan hasil bagi \((u/v)' = (u'v - uv')/v^2\): \[ \delta E = \frac{\delta \langle \psi | \hat{H} | \psi \rangle \cdot \langle \psi | \psi \rangle - \langle \psi | \hat{H} | \psi \rangle \cdot \delta \langle \psi | \psi \rangle}{(\langle \psi | \psi \rangle)^2} = 0 \] Agar hasil bagi ini nol, maka pembilangnya harus nol: \[ \delta \langle \psi | \hat{H} | \psi \rangle \cdot \langle \psi | \psi \rangle = \langle \psi | \hat{H} | \psi \rangle \cdot \delta \langle \psi | \psi \rangle \qquad (3) \]
\end{quote}

\begin{center}\rule{0.5\linewidth}{0.5pt}\end{center}

\hypertarget{jembatan-logika-menuju-persamaan-nilai-eigen}{%
\subsection{Jembatan Logika: Menuju Persamaan Nilai Eigen}\label{jembatan-logika-menuju-persamaan-nilai-eigen}}

Jika kita uraikan variasi pada persamaan (3) terhadap bra \(\langle \psi |\): \[ \langle \delta \psi | \hat{H} | \psi \rangle \cdot \langle \psi | \psi \rangle = \langle \psi | \hat{H} | \psi \rangle \cdot \langle \delta \psi | \psi \rangle \] Pindahkan suku \(\langle \psi | \psi \rangle\) ke bawah: \[ \langle \delta \psi | \hat{H} | \psi \rangle = \left( \frac{\langle \psi | \hat{H} | \psi \rangle}{\langle \psi | \psi \rangle} \right) \langle \delta \psi | \psi \rangle \qquad (4) \]

Karena suku dalam kurung adalah \(E\), maka: \[ \langle \delta \psi | ( \hat{H} | \psi \rangle - E | \psi \rangle ) = 0 \] Karena variasi \(\langle \delta \psi |\) bersifat sembarang, maka bagian dalam kurung harus nol: \[ \hat{H} | \psi \rangle = E | \psi \rangle \qquad (5) \]

\textbf{Kesimpulan Jembatan:} Mencari titik stasioner dari fungsional energi \(E(\psi)\) identik dengan menyelesaikan persamaan nilai eigen Schrodinger.

\begin{center}\rule{0.5\linewidth}{0.5pt}\end{center}

\hypertarget{reduksionisme-pembuktian-epsi-ge-e_0}{%
\subsection{\texorpdfstring{Reduksionisme: Pembuktian \(E(\psi) \ge E_0\)}{Reduksionisme: Pembuktian E(\textbackslash psi) \textbackslash ge E\_0}}\label{reduksionisme-pembuktian-epsi-ge-e_0}}

Sekarang kita buktikan bahwa tebakan apa pun (\(|\psi\rangle\)) selalu memberikan energi di atas atau sama dengan \emph{ground state} (\(E_0\)).

\hypertarget{a.-ekspansi-dalam-basis-eigen}{%
\subsubsection{Ekspansi dalam Basis Eigen}\label{a.-ekspansi-dalam-basis-eigen}}

Misalkan \(\{|v_n\rangle\}\) adalah set vektor eigen dari \(\hat{H}\) dengan energi \(E_0 \le E_1 \le E_2 \dots\). Sembarang state \(|\psi\rangle\) dapat ditulis sebagai: \[ |\psi\rangle = \sum_n c_n |v_n\rangle \qquad (6) \]

\begin{quote}
\textbf{Visualisasi (6): Operasi Matriks} Secara linear, \(|\psi\rangle\) adalah jumlahan berbobot dari state-state murni (basis eigen): \[ |\psi\rangle = c_0 |v_0\rangle + c_1 |v_1\rangle + \dots = \begin{pmatrix} c_0 \\ c_1 \\ \vdots \end{pmatrix} \]
\end{quote}

\hypertarget{b.-substitusi-ke-rayleigh-quotient}{%
\subsubsection{Substitusi ke Rayleigh Quotient}\label{b.-substitusi-ke-rayleigh-quotient}}

\[ E(\psi) = \frac{\sum_n |c_n|^2 E_n}{\sum_n |c_n|^2} \qquad (7) \]

\begin{quote}
\textbf{Visualisasi (7): Perhitungan Linear} Karena \(E_n \ge E_0\) untuk semua \(n\), maka: \[ \sum_n |c_n|^2 E_n \ge \sum_n |c_n|^2 E_0 = E_0 \sum_n |c_n|^2 \] Masukkan kembali ke pembilang persamaan (7): \[ E(\psi) \ge \frac{E_0 \sum_n |c_n|^2}{\sum_n |c_n|^2} = E_0 \qquad (8) \]
\end{quote}

\begin{center}\rule{0.5\linewidth}{0.5pt}\end{center}

\hypertarget{verifikasi-parameter-batas-bawah}{%
\subsection{Verifikasi \& Parameter: Batas Bawah}\label{verifikasi-parameter-batas-bawah}}

\begin{longtable}[]{@{}
  >{\raggedright\arraybackslash}p{(\columnwidth - 4\tabcolsep) * \real{0.3366}}
  >{\raggedright\arraybackslash}p{(\columnwidth - 4\tabcolsep) * \real{0.1683}}
  >{\raggedright\arraybackslash}p{(\columnwidth - 4\tabcolsep) * \real{0.4950}}@{}}
\toprule\noalign{}
\begin{minipage}[b]{\linewidth}\raggedright
Kondisi \(\psi\rangle\)
\end{minipage} & \begin{minipage}[b]{\linewidth}\raggedright
Nilai \(E(\psi)\)
\end{minipage} & \begin{minipage}[b]{\linewidth}\raggedright
Interpretasi Fisik
\end{minipage} \\
\midrule\noalign{}
\endhead
\bottomrule\noalign{}
\endlastfoot
\textbf{\(\psi\rangle =v_0\rangle\)} & \(E(\psi) = E_0\) & Tebakan tepat mengenai \emph{ground state}. \\
\textbf{\(\psi\rangle \perp v_0\rangle\)} & \(E(\psi) \ge E_1\) & Tebakan tidak mengandung komponen \emph{ground state}. \\
\textbf{Sembarang \(\psi\rangle\)} & \(E(\psi) > E_0\) & Tebakan mengandung ``pengotor'' dari state eksitasi. \\
\end{longtable}

\begin{center}\rule{0.5\linewidth}{0.5pt}\end{center}

\hypertarget{physical-insight-natures-compass}{%
\subsection{Physical Insight: ``Nature's Compass''}\label{physical-insight-natures-compass}}

Prinsip variasi bertindak sebagai \textbf{kompas} bagi algoritma VQE. 1. Kita membuat sirkuit parametrik \(|\psi(\theta)\rangle\). 2. Kita menghitung energi \(E(\theta)\) menggunakan QPU. 3. Kita tahu dengan pasti bahwa semakin kecil \(E(\theta)\), semakin dekat kita dengan solusi ekonomi optimal (\(E_0\)).

\textbf{Mengapa ini penting untuk portofolio?} Karena kita tidak bisa mendiagonalkan matriks kovarians yang sangat besar, kita menggunakan Prinsip Variasi untuk ``meraba'' dasar lembah energi. Kita tidak perlu tahu di mana dasarnya, kita hanya perlu terus bergerak ``ke bawah'' sesuai petunjuk Teorema Rayleigh-Ritz.
